\documentclass[conference]{IEEEtran}
\IEEEoverridecommandlockouts
% The preceding line is only needed to identify funding in the first footnote. If that is unneeded, please comment it out.
\usepackage{cite}
\usepackage{amsmath,amssymb,amsfonts}
\usepackage{algorithmic}
\usepackage{graphicx}
\usepackage{placeins}
\usepackage{textcomp}
\usepackage{xcolor}
\usepackage{hyperref}
\def\BibTeX{{\rm B\kern-.05em{\sc i\kern-.025em b}\kern-.08em
    T\kern-.1667em\lower.7ex\hbox{E}\kern-.125emX}}
% Break the IEEEtran author alignment into a new row (useful for a 2x2 layout).
\makeatletter
\newcommand{\linebreakand}{%
  \end{@IEEEauthorhalign}
  \hfill\mbox{}\par
  \mbox{}\hfill\begin{@IEEEauthorhalign}}
\makeatother
\begin{document}

\title{Análise de conjunto de dados de indicadores de diabetes\\
{\footnotesize \textsuperscript{*}Estudo e análise exploratória de um
 conjunto de dados de indicadores de saúde para prever casos de diabetes}
}

\author{\IEEEauthorblockN{1\textsuperscript{st} Filipe Alcantara da Costa}
\IEEEauthorblockA{\textit{dept. Engenharia de Teleinformática (DETI)} \\
\textit{Universidade Federal do Ceará (UFC)}\\
Fortaleza, Brasil  \\
email address or ORCID}
\and
\IEEEauthorblockN{2\textsuperscript{nd} Gustavo Oliveira Seabra}
\IEEEauthorblockA{\textit{dept. Engenharia de Teleinformática (DETI)} \\
\textit{Universidade Federal do Ceará (UFC)}\\
Fortaleza, Brasil  \\
gustavooliveiraseabra@alu.ufc.br}
\linebreakand
\IEEEauthorblockN{3\textsuperscript{rd} Lucas de Oliveira Sobral}
\IEEEauthorblockA{\textit{dept. Engenharia de Teleinformática (DETI)} \\
\textit{Universidade Federal do Ceará (UFC)} \\
Fortaleza, Brasil  \\
sobrallucas33@gmail.com}
\and
\IEEEauthorblockN{4\textsuperscript{th} Savlio Carvalho Pontes}
\IEEEauthorblockA{\textit{dept. Engenharia de Teleinformática (DETI)} \\
\textit{Universidade Federal do Ceará (UFC)}\\
Fortaleza, Brasil \\
savliocp@alu.ufc.br}
}

\maketitle

\begin{abstract}
Neste primeiro trabalho da cadeira de Inteligência Computacional Aplicada, realizamos a primeira etapa para nosso futuro modelo Preditivo ou de Inferência, sendo esta, a etapa de análise exploratória do dados, é a etapa mais importante e que leva maior tempo, seja no Aprendizado estatístico ou Aprendizado de Máquina.

Aqui estabelecemos o escopo e base para nosso futuro modelo, suas variáveis predioras, seu objetivo e variável-alvo, verificamos erros da base de dados, transformamos seus preditores.

Utilizamos para isso várias etapas de exploração, sendo elas: Pré-processamento de dados, Análise Univariada Condicional e Não-Condicional, Análise Bivariada Condicional e Não-Condicional, Análise Multivariada através da criação dos Componentes Principais (PCA).

Nossos resultados foram a definição de 6 preditores definitivos dentre os 22 totais, uma variável-alvo representada pela classe \texttt{Diabetes\_binary} e duas componentes principais que explicam 59,41\% dos dados, com o intuito de, futuramente, desenvolver um modelo capaz de prever e inferir se uma pessoa tem diabetes.

\end{abstract}

\begin{IEEEkeywords}
Análise Exploratória, Pré-processamento, PCA, Diabetes, Análise de dados
\end{IEEEkeywords}

\section{Introdução}

A diabetes é uma condição caracterizada pelo aumento do nível de glicose no sangue. O que ocorre quando o corpo de uma pessoa não produz suficientemente ou não consegue usar o hormônio insulina, responsável por levar a glicose às células.
Essa condição pode ser tanto adquirida por influências genéticas, como também adquirida ao longo da vida, ligada a fatores como o sedentarismo, hábitos alimentares e sobrepeso.
Em ambos os casos, a diabetes tem uma forte influência na vida da pessoa, ocasionando sintomas como fadiga, visão turva, perda de peso, feridas, dormência, sede e infecções.

No Brasil, essa condição é bem prevalente. Segundo o Ministério da Saúde, a frequência do diagnóstico foi de 10,2\%, o que corresponde a cerca de 20 milhões da população, ocupando o sexto lugar entre países com mais pessoas com diabetes.

Em meio a esse cenário, este estudo faz uso de um dataset fornecido pela plataforma Kaggle: 
\href{https://www.kaggle.com/datasets/alexteboul/diabetes-health-indicators-dataset/data?select=diabetes_binary_5050split_health_indicators_BRFSS2015.csv}{Conjunto de dados de indicadores de saúde para diabetes},
com o intuito de fazer uma Análise Exploratória de Dados (EDA) para entender as características e hábitos de uma pessoa e sua relação com a diabetes.

Dos preditores do dataset, dois serão de maior relevância para as nossas análises: \textit{Diabetes\_binary}, representando se a pessoa tem diabetes ou não, e \textit{BMI}, 
índice de massa corporal da pessoa. Essas serão nossas variáveis-alvo, as quais nos ajudarão a representar se os demais preditores são ou não úteis para a nossa análise.

Essa análise exploratória dos dados é importante para o entendimento do dataset como um todo. Identificando as nuances de cada preditor e a relação deles com os outros, o que ajudará em outros estudos para realizar modelos de regressão ou classificação. Isso por meio da redução da dimensionalidade do dataset, como também pela identificação de outliers e pela análise das relações de preditores com nossas variáveis-alvo.



\section{Metodologia}

Trata-se de uma análise de objetivo descritivo para entender a relação entre a diabetes e os demais preditores, utilizando, para isso, abordagem quantitativa, por meio da manipulação numérica do dataset. O estudo também foca na análise de caso para investigar e responder
à pergunta central de qual é a relação da diabetes com outras características de uma pessoa, e qual característica está mais relacionada a essa condição,
sendo, de natureza, uma pesquisa básica.


O estudo foi realizado usando o dataset do Kaggle referido integrado ao ambiente do Google Colab, para a utilização de Python.
Foi utilizado das bibliotecas kagglehub, para a importação dos dados, numpy, pandas e scipy para a manipulação do dataset, além da seaborn e da matplotlib 
para o plot dos gráficos. Os gráficos foram criados com base no tipo de variável que estava sendo analisada em questão.

Então, a análise foi dividida em 4 partes: o entendimento e tratamento dos dados, a análise univariada, a análise bivariada e a análise
multivariada.


\subsection{Primeira análise dos dados}

O primeiro passo foi importar o dataset do \href{https://www.kaggle.com/datasets/alexteboul/diabetes-health-indicators-dataset/data?select=diabetes_binary_5050split_health_indicators_BRFSS2015.csv}{Kaggle}. 
Esse dataset conta, a priori, com 70.692 amostras e 22 variáveis. Esse dataset é um dataset já identado de outro usado para uma pesquisa, logo 
ele não apresenta valores faltando e a quantidade de pessoas com diabetes é igual à de pessoas sem diabetes. Entretanto, ainda assim, 
irá conter amostras com valores duplicados e outliers.

Para evitar trabalhar com uma quantidade massiva de dados, visto que dificultaria a visualização, foi então feito um processo de tratamento.
Primeiro foi feita uma randomização, para tirar o viés, e coleta de uma parte dos dados, com 6.000 amostras restantes.
Após isso, foram retiradas amostras duplicadas.

Depois desse processo de tratamento, foi retirado o dataset de 5.983 amostras e 22 preditores. 
Os preditores constam com 3 tipos no total:

\begin{itemize}
    \item 15 variáveis binárias, sendo elas: "Diabetes\_binary, HighBP, HighChol, CholCheck, Smoker, Stroke, HeartDiseaseorAttack, 
physActivity, Fruits, Veggies, HvyAlcoholConsump, AnyHealthcare, NoDocbcCost, DiffWalk e Sex";
    \item 4 qualitativas ordinais: "GenHlth, Age, Education e Income";
    \item 3 quantitativas: "MentHlth, PhysHlth, BMI (Body Mass Index)".
\end{itemize}
 

Um objetivo para artigos posteriores a este é fazer um modelo de regressão e um de classificação. Para tanto, será escolhida uma 
variável principal de análise para cada modelo, as quais, desde já, serão tratadas como variáveis-alvo para o foco da análise.
Para o modelo de regressão, precisamos de uma variável quantitativa; para tanto, o BMI foi escolhido. Já para a classificação, 
como o nosso objetivo é comparar a diabetes com os demais preditores, foi escolhido o preditor "Diabetes\_binary".

A partir dessa escolha, o preditor "Diabetes\_binary" também foi escolhido para a separação por classes, o que será útil para a análise univariada dos preditores, sendo no total 2 classes: a 0, sem diabetes, e a 1, com diabetes. 

O nosso objetivo com as análises é:
\begin{itemize}
    \item Diminuir a dimensionalidade do dataset;
    \item Verificar a relação dos preditores com as variáveis-alvo;
    \item Verificar a relação dos preditores entre si.
\end{itemize}

Para tanto, em cada análise, serão mostrados gráficos para contribuir com esses objetivos.


\subsection{Analise Univariada}

Na análise Univariada, o nosso objetivo é analisar o tipo e o comportamento do preditor, além de comparar as classes de um mesmo preditor 
para verificar a relação com diabets. Procura-se preditores que tenham mais diferenças em casos com e sem diabets, para termos um aspecto de quais 
variaveis são mais relevantes para os modelos. Para tanto, foi feito tanto um plot das medidas de tendencias gerais, média, desvio 
padrão e assimetria, como tambem de graficos para a comparação entre classes, de acordo com o tipo de variavel.

Os graficos para a visualização, se relacionam muito com o tipo da variavel:

\begin{itemize}
    \item Binarios: porcentagem de respostas afirmativas em relação ao total;
    \item Qualitativos: graficos de linhas ou densidade.
    \item Quantitativos: histogramas e boxplots;
\end{itemize}

\subsubsection{Variaveis Binárias} 

Para a análise univariada dos preditores Binários, foi usado o grafico da 
figura \ref{fig:Binários}. Nele, é possivel observar quantas amostras tem o valor 1
em relação ao total do preditor, em porcentagem. Em azul, está o dataset normal, já em
laranja e verde, está a separação por classes de acordo com o preditor "Diabets\_Binary".

Esse grafico consegue fornecer uma interpretação muito boa de um preditor 
conseguimos ver a tendencia, de forma que, quanto mais perto dos extremos, 
mais respostas iguais para um preditor irão ser comuns, já em relação a
diabets, quanto mais proximo as duas cores estão, em seus maximos, menos relação
aquele preditor terá com a diabets, ao passo que, quanto maior a variação, mais provavelmente
o preditor estará relacionado à diabets. Como "diabets\_binary" é a variavel alvo, o foco estará em
preditores que estejam mais ao centro e com mais variação entre as classes.

De acordo com esses dois critérios, foi possivel focar em alguns preditores mais relevantes
para o caso, sendo eles:

\begin{itemize}
    \item \textbf{HighChol}: se a pessoa tem coresterol alto;
    \item \textbf{HighBP}: se a pessoa tem pressão alta;
    \item \textbf{physActivity}: se a pessoa faz atividade física
    \item \textbf{DiffWalk}:  se tem dificultade para andar em escadas;
    \item \textbf{HeartDiseaseorAttack}: se tem doença arterial ou já teve infarto;
\end{itemize}





\begin{figure*}[!t]
    \centering
    
    % 1. FIGURA DE CIMA (Tela toda)
    \includegraphics[width=1\textwidth]{Figures/univariavel_binarias.png}
    \caption{Quantidade de respostas afirmativas em relação ao total das varivaveis binárias}

    \label{fig:Binários}
\end{figure*}  
\FloatBarrier

\subsubsection{Variaveis qualitativas}

Para a análise das variaveis qualitativas, como possuem mais valores que somente dois, é necessário um grafico mais detalhado.
Para tanto, foi escolhida a ultilização de um grafico de linhas representando a porcentagem de cada categoria em relação ao total, assim como 
no grafico da figura \ref{fig:GentHlth}. Esse preditor, GentHlth, ou, nivel de saúde, possui 5 classificações, indo de 1 a 5, sendo 1 excelente e 5 muito ruim.
No grafico, é possivel ver a distribuição de cada classe, tanto para o dataset normal, em azul, como para as classes, laranja para sem diabets
e verde para com diabets. 

Esse exemplo seve para mostrar como analisar esse tipo de grafico. É possivel ver um agrupamento de mais pessoas diabeticas em niveis altos, enquanto pessoas não diabeticas
Estão agrupados para niveis baixos. Essa relação mostra como essa muda de acordo com o preditor "Diabets\_Binary", sendo um grande indicativo que será útil para as análises.

O mesmo grafico foi feito para os demais preditores qualitativos, Age, Education e Income. Os preditores Age e GentHlth foram
os que apresentaram maior variação, enquanto os preditores Income e Education, não mostraram tanta variação, porém ainda sendo relevantes.

Os preditores são:
\begin{itemize}
    \item \textbf{GentHlth} Nível de sauúde, de 1 a 5, sendo 1 excelente e 5 pessima;
    \item \textbf{Age}: Categoria etária de 13 níveis segundo "\_AGEG5YR";
    \item \textbf{Education}: Nível de escolaridade, de 1 a 6; 
    \item \textbf{Income}: nível de renda, de 1 a 8.
\end{itemize}

\begin{figure}[!ht]
    \centering
    \includegraphics[width=0.7\columnwidth]{Figures/Genhlth.png}
    \caption{Grafico de linhas, em porcentagem, do preditor GentHlth.}
    \label{fig:GentHlth}
\end{figure}
\FloatBarrier

\subsubsection{Variaveis quantitatvas}

Para a análise das variaveis quantitativas, como elas possuem mais valores possiveis,
foi decidido fazer um histograma em conjunto para a analise de distribuição.
O exemplo pode ser visto na figura \ref{fig:BMI}, que representa o histrograma com a linha de densidade
do preditor BMI. Em azul, está o dataset normal, em amrelo a classe sem diabetes e em verde, a calasse com diabets.

Por esses graficos é possi notar tanto a tendencia de concentração, com um agrupamento em aproximadamente 30,
como também a diferençaentre as classes. 
a classe com diabets aprece estar deslocada para a direira, o que induz a pensar 
em uma relação leve entre os preditores, diforma que pessoas com diabetes tem mais BMI.

Também foram feitos graficos similares para PhysHlth e MentHlth. Ambos tem uma concentração maior em 0, 
o PhysHlth parece ter um grande deslocamento em relação as classes com e sem diabets.
Então, todos os preditores parecem uteis para as análises.

Os preditores são:
\begin{itemize}
    \item \textbf{BMI}: indice de massa corporal;
    \item \textbf{PhysHlth}: quantidade de dias com saúde física debilitada nos ultimos 30 dias;
    \item \textbf{MentHlth}:  quantidade de dias com saúde mental debilitada nos ultimos 30 dias.
\end{itemize}

\begin{figure}[!ht]
    \centering
    \includegraphics[width=0.7\columnwidth]{Figures/BMI.png}
    \caption{histograma do preditor BMI, com e sem classificação.}
    \label{fig:BMI}
\end{figure}
\FloatBarrier


\subsection{Análise Bivariada}

\subsection{Análise Multivariada}
A Análise Multivariada neste trabalho implementou e utilizou a Análise de Componentes Principais (PCA), que é uma técnica estatística multivariada e de aprendizado não supervisionado, voltada para a redução de dimensionalidade e extração de características (feature extraction). Como é um método não supervisionado, ele não utiliza a variável resposta $Y$ ou rótulos de classe para construir seus componentes. 

A sua primeira aplicação seguiu os seguintes passos:

\begin{enumerate}
  \item \textbf{Obter matriz de Covariância ou Correlação}: Levando em conta nosso dataset, que havia tanto dados contínuos quanto categóricos/binários, optamos por usar a Matriz de Correlação, pois ela é mais indicada para preditores que estão em escalas diferentes ou unidades de medidas diferentes, em oposição à de Covariância.

  \item \textbf{Decomposição em Autovalores e Autovetores}: Após se obter a Matriz de Correlação, podemos calcular seus autovetores, que definem a direção e a orientação dos novos eixos ortogonais no espaço multidimensional, e seus autovalores, que representam a magnitude da variância explicada ao longo da direção do seu autovetor correspondente; por fim, deixá-los em ordem decrescente para que o maior e mais importante componente fique em primeiro na ordem.

  \item \textbf{Cálculo da variância explicada}: Se trata do cálculo dos autovalores em ordem decrescente pela razão do somatório de todos os autovalores, multiplicado por cem para termos o cálculo em porcentagem, sendo o primeiro componente ($PC_1$) a combinação linear normalizada que maximiza a variância dos dados e o segundo componente principal ($PC_2$) o que captura a segunda maior fatia da variância dos dados, sob restrição de ser ortogonal (com correlação 0) ao $PC_1$.

  \item \textbf{Projeção dos Dados (Factor Scores)}: E por fim podemos plotar o gráfico de calor (Heatmap) com os dois principais componentes ou todos os componentes principais, juntamente com o gráfico Biplot do $PC_1$ e $PC_2$, sendo justamente uma das vantagens da análise de componentes principais, pois não podíamos plotar um gráfico 2D com a nossa quantidade de preditores, mas com os dois principais componentes que explicam o conjunto de dados podemos observar seu comportamento vetorial.

\end{enumerate}

Em sua segunda aplicação:

\begin{enumerate}
  \item \textbf{Escolha dos preditores}: A primeira aplicação do PCA explicou apenas $24,54\%$ da variância acumulada. Além disso, a variável-alvo $Diabetes\_binary$ havia sido incluída entre os preditores, o que causaria vazamento de dados no treinamento e no teste. Por isso, ela foi removida do cálculo.
  
  \item \textbf{Remoção de Preditores por Multicolinearidade}: Sem a variável-alvo, a variância explicada aumentou para $34,60\%$. Avaliamos então a remoção de preditores redundantes por multicolinearidade, mas a maior correlação encontrada foi de apenas $0,51$, entre $DiffWalk$ e $PhysHlth$. Assim, essa abordagem não permitiria uma redução significativa.

  \item \textbf{Combinação dos melhores preditores}: Como a baixa correlação impedia a seleção por redundância, aplicamos um algoritmo combinatório para avaliar subconjuntos com 4, 5, 6 e 7 preditores. Os melhores resultados são apresentados na Tabela \ref{tab:melhores_combinacoes_pca}.

  \item \textbf{Preditores definitivos}: A combinação de 6 preditores apresentou o melhor equilíbrio entre dimensionalidade e variância explicada acumulada ($59,41\%$): HighBP, HighChol, MentHlth, PhysHlth, Age e GenHlth.


\end{enumerate}

\begin{table*}[!t]
  \centering
  \caption{Melhores combinações de variáveis segundo a variância explicada por PC1 e PC2.}
  \label{tab:melhores_combinacoes_pca}
  \begin{tabular}{c c l}
  \hline
  \textbf{Qtd. de variáveis} &
  \textbf{Variância explicada (\%)} &
  \textbf{Variáveis} \\
  \hline

  4 & 76,67 &
  PhysHlth, Income, Education, GenHlth \\

  5 & 67,73 &
  MentHlth, PhysHlth, Income, Education, GenHlth \\

  6 & 59,41 &
  HighBP, HigChol, MentHlth, PhysHlth, Age, Income \\

  7 & 54,20 &
  HighBP, HighChol, MentHlth, PhysHlth, Age, ... \\

  \hline
  \end{tabular}
\end{table*}
\FloatBarrier


\subsection{Estrutura dos dados e checagens iniciais}
O processo de amostragem aleatória e limpeza resultou em um conjunto final de 5.983 amostras e 22 preditores, livre de dados nulos. As checagens iniciais confirmaram o equilíbrio das classes na variável-alvo \textit{Diabetes\_binary} (proporção idêntica de indivíduos com e sem diabetes) e a presença de \textit{outliers} e assimetria em variáveis contínuas, como o \textit{BMI}, \textit{PhysHlth} e \textit{MentHlth}. Essas características reforçaram a necessidade de padronização prévia dos dados para a aplicação do PCA.

\subsection{Achados principais}

A análise univariada evidenciou comportamentos distintos entre os indivíduos com e sem diabetes ao longo das três categorias de preditores:

\begin{itemize}
    \item \textbf{Preditores Binários}: A comparação percentual mostrou que o grupo com diabetes (classe 1) apresenta prevalências substancialmente maiores em diagnósticos clínicos e limitações físicas quando comparado ao grupo sem diabetes (classe 0). Pode-se observar esse comportamento em \textit{HighBP} ($74,8\%$  x $37,2\%$ ), \textit{HighChol} ($63,7\%$ x $37,8\%$), \textit{DiffWalk} ($36,9\%$ x $13,2\%$) e \textit{HeartDiseaseorAttack} ($18,6\%$ x $5,2\%$). Em contrapartida, a prática de atividade física (\textit{PhysActivity}) foi menor entre os diabéticos ($63,1\%$ x $77,9\%$).
    
    \item \textbf{Preditores Qualitativos}: O preditor \textit{GenHlth} (autoavaliação de saúde) demonstrou forte capacidade discriminatória: a classe sem diabetes concentrou-se nos níveis 1 e 2 (saúde excelente/boa), enquanto a classe com diabetes apresentou distribuição deslocada para os níveis 3, 4 e 5 (saúde razoável/ruim/pessima). A variável \textit{Age} também confirmou que a prevalência de diabetes cresce de forma contínua com o avanço das faixas etárias. Os preditores socioeconômicos \textit{Income} e \textit{Education} indicaram leve tendência de menor renda e menor escolaridade no grupo afetado pela condição.

    \item \textbf{Preditores Quantitativos}: O \textit{BMI} exibiu média superior na classe com diabetes ($31,87 \text{ kg/m}^2$, desvio padrão $7,32$) em relação à classe sem diabetes ($27,78 \text{ kg/m}^2$, desvio padrão $5,82$), com o histograma de densidade demonstrando um deslocamento nítido da curva à direita. Para \textit{PhysHlth} e \textit{MentHlth}, embora a mediana de ambas as classes seja zero, a média de dias com saúde física debilitada foi maior nos diabéticos ($7,90$ dias vs. $2,83$ dias nos últimos 30 dias).
\end{itemize}

\subsection{Redundância e seleção}
A avaliação de associação bivariada via coeficiente de correlação de Pearson revelou que, no geral, os preditores do conjunto de dados apresentam baixas correlações lineares interpessoais. A maior correlação observada ocorreu entre as variáveis \textit{PhysHlth} e \textit{DiffWalk} ($r = 0,51$), seguida por \textit{GenHlth} e \textit{PhysHlth} ($r = 0,50$). 

A aplicação do algoritmo de filtragem por limiar de multicolinearidade confirmou a ausência de redundâncias extremas ($r > 0,80$), o que inviabilizou a eliminação direta de atributos por simples correlação de pares. Isso justificou o uso da abordagem combinatória com PCA para a escolha do melhor subconjunto reduzido.

\subsection{Análise dos Componentes Principais (PCA) e Biplot}


\section{Conclusão}



\begin{thebibliography}{00}
\bibitem{b1} G. Eason, B. Noble, and I. N. Sneddon, ``On certain integrals of Lipschitz-Hankel type involving products of Bessel functions,'' Phil. Trans. Roy. Soc. London, vol. A247, pp. 529--551, April 1955.
\bibitem{b2} J. Clerk Maxwell, A Treatise on Electricity and Magnetism, 3rd ed., vol. 2. Oxford: Clarendon, 1892, pp.68--73.
\bibitem{b3} I. S. Jacobs and C. P. Bean, ``Fine particles, thin films and exchange anisotropy,'' in Magnetism, vol. III, G. T. Rado and H. Suhl, Eds. New York: Academic, 1963, pp. 271--350.
\bibitem{b4} K. Elissa, ``Title of paper if known,'' unpublished.
\bibitem{b5} R. Nicole, ``Title of paper with only first word capitalized,'' J. Name Stand. Abbrev., in press.
\bibitem{b6} Y. Yorozu, M. Hirano, K. Oka, and Y. Tagawa, ``Electron spectroscopy studies on magneto-optical media and plastic substrate interface,'' IEEE Transl. J. Magn. Japan, vol. 2, pp. 740--741, August 1987 [Digests 9th Annual Conf. Magnetics Japan, p. 301, 1982].
\bibitem{b7} M. Young, The Technical Writer's Handbook. Mill Valley, CA: University Science, 1989.
\end{thebibliography}
\vspace{12pt}

\end{document}